\subsection{Duality and Multilinear Structure}

% **Role:** algebraic language of linear structure.
%
% Dual spaces, dual maps, annihilators, transpose as dual map, bilinear forms, tensor products as later extension.

Duality studies vectors through scalar-valued tests. Functionals, dual maps, annihilators, transposes, bilinear forms, and tensor products extend linear structure into multilinear language.

\begin{definition} \hlt{Linear Functional and Dual Space}\\
A \hlt{linear functional} on $\mathcal{V}$ is a linear map from $\mathcal{V}$ to $\mathbb{F}$. The \hlt{dual space} of $\mathcal{V}$ is $\mathcal{V}' \defeq \LL(\mathcal{V},\mathbb{F})$.
\end{definition}

\begin{definition} \hlt{Dual Basis}\\
Let $\beta=(\mathbf{v}_1,\ldots,\mathbf{v}_n)$ be a basis of $\mathcal{V}$. The \hlt{dual basis} of $\beta$ is the list $\beta'=(\varphi_1,\ldots,\varphi_n)$ in $\mathcal{V}'$ defined by
\begin{align}
\varphi_j(\mathbf{v}_k)=
\begin{cases}
1, & k=j,\\
0, & k\neq j.
\end{cases}
\nonumber
\end{align}
Existence and uniqueness of each $\varphi_j$ follow from Proposition \ref{prop:linear-map-determined-by-basis}.
\end{definition}

\begin{theorem} \label{thm:dual-basis-and-dimension} \hlt{Dual Basis and Dimension of the Dual Space}\\
Let $\mathcal{V}$ be finite-dimensional with basis $\beta=(\mathbf{v}_1,\ldots,\mathbf{v}_n)$, and let $\beta'=(\varphi_1,\ldots,\varphi_n)$ be the dual basis. Then $\beta'$ is a basis of $\mathcal{V}'$. In particular, $\dim(\mathcal{V}')=\dim(\mathcal{V})$.
\begin{proof}
First suppose $a_1\varphi_1+\cdots+a_n\varphi_n=0$. Evaluating at $\mathbf{v}_j$ gives $a_j=0$ for each $j$, so $\beta'$ is linearly independent. Now let $\varphi\in\mathcal{V}'$. For $\mathbf{v}=c_1\mathbf{v}_1+\cdots+c_n\mathbf{v}_n$, $\left(\sum\limits_{j=1}^n \varphi(\mathbf{v}_j)\varphi_j\right)(\mathbf{v}) =\sum\limits_{j=1}^n \varphi(\mathbf{v}_j)c_j =\varphi(\mathbf{v})$.\\
Thus $\varphi=\sum_{j=1}^n \varphi(\mathbf{v}_j)\varphi_j$, so $\beta'$ spans $\mathcal{V}'$.\\
\hlt{Structural Commentary}: A basis of $\mathcal{V}$ creates coordinate-reading functionals on $\mathcal{V}$. These functionals form the coordinate system for the dual space.\\
\hlt{Key Insight}: A linear functional is determined exactly by its values on a basis.
\end{proof}
\end{theorem}

\begin{definition} \label{def:dual-map} \hlt{Dual Map}\\
If $T\in\LL(\mathcal{V},\mathcal{W})$, the \hlt{dual map} of $T$ is the map $T':\mathcal{W}'\to\mathcal{V}'$ defined by $T'\varphi \defeq \varphi T$ for every $\varphi\in\mathcal{W}'$.\\
Equivalently, $(T'\varphi)(\mathbf{v})=\varphi(T\mathbf{v})$.
\end{definition}

\begin{proposition} \label{prop:dual-map-is-linear} \hlt{Dual Maps Are Linear}\\
If $T\in\LL(\mathcal{V},\mathcal{W})$, then $T'\in\LL(\mathcal{W}',\mathcal{V}')$.
\begin{proof}
For $\varphi,\psi\in\mathcal{W}'$, $a\in\mathbb{F}$, and $\mathbf{v}\in\mathcal{V}$, $\left(T'(a\varphi+\psi)\right)(\mathbf{v}) =(a\varphi+\psi)(T\mathbf{v}) =a\varphi(T\mathbf{v})+\psi(T\mathbf{v}) =\left(aT'\varphi+T'\psi\right)(\mathbf{v})$.\\
\hlt{Structural Commentary}: The dual map sends a test on $\mathcal{W}$ to the pulled-back test on $\mathcal{V}$. Linearity is inherited from pointwise linearity of functionals.\\
\hlt{Key Insight}: Dual maps reverse direction because a functional can test $T\mathbf{v}$ only after $T$ is applied.
\end{proof}
\end{proposition}

\begin{theorem} \label{thm:algebra-of-dual-maps} \hlt{Algebra of Dual Maps}\\
The dual operation has the following properties:
\begin{enumerate}[label=(\roman*)]
\setlength{\itemsep}{0pt}
\item If $S,T\in\LL(\mathcal{V},\mathcal{W})$ and $a\in\mathbb{F}$, then $(aS+T)'=aS'+T'$.
\item If $T\in\LL(\mathcal{U},\mathcal{V})$ and $S\in\LL(\mathcal{V},\mathcal{W})$, then $(ST)'=T'S'$.
\item $(I_{\mathcal{V}})'=I_{\mathcal{V}'}$.
\end{enumerate}
\begin{proof}
\begin{enumerate}[label=(\roman*)]
\setlength{\itemsep}{0pt}
\item For $\varphi\in\mathcal{W}'$ and $\mathbf{v}\in\mathcal{V}$, $\left((aS+T)'\varphi\right)(\mathbf{v}) =\varphi(aS\mathbf{v}+T\mathbf{v}) =a(S'\varphi)(\mathbf{v})+(T'\varphi)(\mathbf{v})$, so $(aS+T)'=aS'+T'$.
\item For $\varphi\in\mathcal{W}'$ and $\mathbf{u}\in\mathcal{U}$, $\left((ST)'\varphi\right)(\mathbf{u}) =\varphi(S(T\mathbf{u})) =\left(T'(S'\varphi)\right)(\mathbf{u})$, so $(ST)'=T'S'$.
\item For $\varphi\in\mathcal{V}'$ and $\mathbf{v}\in\mathcal{V}$, $(I_{\mathcal{V}}'\varphi)(\mathbf{v})=\varphi(I_{\mathcal{V}}\mathbf{v})=\varphi(\mathbf{v})$, so $(I_{\mathcal{V}})'=I_{\mathcal{V}'}$.
\end{enumerate}
\hlt{Structural Commentary}: Duality preserves linear combinations but reverses products. This reversal is the same order reversal seen in transposes.\\
\hlt{Key Insight}: Pullbacks compose in the opposite order.
\end{proof}
\end{theorem}

\begin{definition} \hlt{Annihilator}\\
If $\mathcal{U}\subseteq\mathcal{V}$, the \hlt{annihilator} of $\mathcal{U}$ is $\mathcal{U}^0\defeq\{\varphi\in\mathcal{V}':\varphi(\mathbf{u})=0 \text{ for every } \mathbf{u}\in\mathcal{U}\}$.
\end{definition}

\begin{proposition} \label{prop:annihilator-is-subspace} \hlt{Annihilators Are Subspaces}\\
If $\mathcal{U}\subseteq\mathcal{V}$, then $\mathcal{U}^0$ is a subspace of $\mathcal{V}'$.
\begin{proof}
The zero functional lies in $\mathcal{U}^0$. If $\varphi,\psi\in\mathcal{U}^0$, $a\in\mathbb{F}$, and $\mathbf{u}\in\mathcal{U}$, then $(a\varphi+\psi)(\mathbf{u})=a\varphi(\mathbf{u})+\psi(\mathbf{u})=0$.\\
Thus $a\varphi+\psi\in\mathcal{U}^0$, and Theorem \ref{thm:subspacetest} shows that $\mathcal{U}^0$ is a subspace of $\mathcal{V}'$.\\
\hlt{Structural Commentary}: The annihilator collects all linear tests that vanish on a chosen set. Closure follows because vanishing is preserved by linear combinations of tests.\\
\hlt{Key Insight}: Annihilators are kernels of evaluation constraints.
\end{proof}
\end{proposition}

\begin{theorem} \label{thm:dimension-of-annihilator} \hlt{Dimension of the Annihilator}\\
Let $\mathcal{V}$ be finite-dimensional and let $\mathcal{U}$ be a subspace of $\mathcal{V}$. Then $\dim(\mathcal{U})+\dim(\mathcal{U}^0)=\dim(\mathcal{V})$.
\begin{proof}
Choose a basis $(\mathbf{u}_1,\ldots,\mathbf{u}_m)$ of $\mathcal{U}$ and extend it to a basis $(\mathbf{u}_1,\ldots,\mathbf{u}_m,\mathbf{w}_1,\ldots,\mathbf{w}_k)$ of $\mathcal{V}$.\\
Let $(\varphi_1,\ldots,\varphi_m,\psi_1,\ldots,\psi_k)$ be the corresponding dual basis.\\
By Theorem \ref{thm:dual-basis-and-dimension}, every $\varphi\in\mathcal{V}'$ has the expansion $\varphi=\sum\limits_{j=1}^m \varphi(\mathbf{u}_j)\varphi_j+\sum\limits_{\ell=1}^k \varphi(\mathbf{w}_\ell)\psi_\ell$.\\
If $\varphi\in\mathcal{U}^0$, then $\varphi(\mathbf{u}_j)=0$ for every $j$, so $\varphi\in\operatorname{span}(\psi_1,\ldots,\psi_k)$.\\
Conversely, each $\psi_\ell$ vanishes on every $\mathbf{u}_j$, hence on $\mathcal{U}=\operatorname{span}(\mathbf{u}_1,\ldots,\mathbf{u}_m)$. Therefore $\mathcal{U}^0=\operatorname{span}(\psi_1,\ldots,\psi_k)$.\\
The list $(\psi_1,\ldots,\psi_k)$ is linearly independent as a sublist of the dual basis, so $\dim(\mathcal{U}^0)=k=\dim(\mathcal{V})-\dim(\mathcal{U})$.\\
\hlt{Structural Commentary}: Extending a basis of $\mathcal{U}$ separates directions inside $\mathcal{U}$ from complementary directions. The annihilator is dual to the complementary part.\\
\hlt{Key Insight}: To vanish on $\mathcal{U}$, a functional can only use dual coordinates attached to directions outside $\mathcal{U}$.
\end{proof}
\end{theorem}

\begin{theorem} \label{thm:dual-kernel-range-annihilator} \hlt{Kernel and Range of the Dual Map}\\
Let $\mathcal{V}$ and $\mathcal{W}$ be finite-dimensional, and let $T\in\LL(\mathcal{V},\mathcal{W})$. Then
\begin{enumerate}[label=(\roman*)]
\setlength{\itemsep}{0pt}
\item $\ker(T')=(\operatorname{range}(T))^0$;
\item $\dim(\ker(T'))=\dim(\ker(T))+\dim(\mathcal{W})-\dim(\mathcal{V})$;
\item $\operatorname{range}(T')=(\ker(T))^0$;
\item $\dim(\operatorname{range}(T'))=\dim(\operatorname{range}(T))$.
\end{enumerate}
\begin{proof}
\begin{enumerate}[label=(\roman*)]
\setlength{\itemsep}{0pt}
\item
For $\varphi\in\mathcal{W}'$, $T'\varphi=0$ if and only if $\varphi(T\mathbf{v})=0$ for every $\mathbf{v}\in\mathcal{V}$, which is equivalent to $\varphi\in(\operatorname{range}(T))^0$.
\item By Theorems \ref{thm:dimension-of-annihilator} and \ref{thm:rank-nullity}, $\dim(\ker(T')) =\dim(\mathcal{W})-\dim(\operatorname{range}(T)) =\dim(\mathcal{W})-\dim(\mathcal{V})+\dim(\ker(T))$.
\item If $\varphi=T'\psi$ and $\mathbf{v}\in\ker(T)$, then $\varphi(\mathbf{v})=\psi(T\mathbf{v})=0$, so $\operatorname{range}(T')\subseteq(\ker(T))^0$,\\
Conversely, suppose $\varphi\in(\ker(T))^0$. Define $\widetilde{\varphi}:\operatorname{range}(T)\to\mathbb{F}$ by $\widetilde{\varphi}(T\mathbf{v})=\varphi(\mathbf{v})$.\\
This is well-defined because $T\mathbf{v}=T\mathbf{u}$ implies $\mathbf{v}-\mathbf{u}\in\ker(T)$, hence $\varphi(\mathbf{v}-\mathbf{u})=0$ and $\varphi(\mathbf{v})=\varphi(\mathbf{u})$.\\
It is linear since $\widetilde{\varphi}\big(aT\mathbf{v}+T\mathbf{u}\big) =\widetilde{\varphi}\big(T(a\mathbf{v}+\mathbf{u})\big) =\varphi(a\mathbf{v}+\mathbf{u}) =a\varphi(\mathbf{v})+\varphi(\mathbf{u})$.\\
Choose a basis of $\operatorname{range}(T)$ and extend it to a basis of $\mathcal{W}$. Assign to the basis vectors in $\operatorname{range}(T)$ the values prescribed by $\widetilde{\varphi}$, assign $0$ to the added basis vectors, and use Proposition \ref{prop:linear-map-determined-by-basis} to obtain $\psi\in\mathcal{W}'$ extending $\widetilde{\varphi}$. For every $\mathbf{v}\in\mathcal{V}$, $(T'\psi)(\mathbf{v})=\psi(T\mathbf{v})=\widetilde{\varphi}(T\mathbf{v})=\varphi(\mathbf{v})$, so $\varphi\in\operatorname{range}(T')$. Hence $\operatorname{range}(T')=(\ker(T))^0$.
\item By rank-nullity applied to $T'$ and by (i),\\
$\dim(\operatorname{range}(T')) =\dim(\mathcal{W}')-\dim(\ker(T')) =\dim(\mathcal{W})-\dim((\operatorname{range}(T))^0) =\dim(\operatorname{range}(T))$,\\
where the last equality uses Theorems \ref{thm:dual-basis-and-dimension} and \ref{thm:dimension-of-annihilator}.
\end{enumerate}
\hlt{Structural Commentary}: Dual kernels are annihilators of ranges, and dual ranges are annihilators of kernels. This is the rank-nullity geometry of $T$ seen through scalar tests.\\
\hlt{Key Insight}: A functional dies after pulling back by $T$ exactly when it vanishes on everything $T$ can reach.
\end{proof}
\end{theorem}

% Source note: the prompt-supplied linear-programming screenshots state the equality-only alternative as $A\mathbf{x}=\mathbf{b}$ versus a row-vector certificate. The cone and inequality versions are not used here.
\begin{theorem} \label{thm:range-membership-dual-alternative} \hlt{Range Membership Dual Alternative}\\
Let $\mathcal{V},\mathcal{W}$ be finite-dimensional, $T\in\LL(\mathcal{V},\mathcal{W})$, and $\mathbf{w}\in\mathcal{W}$. Exactly one holds:
\begin{enumerate}[label=(\roman*)]
\setlength{\itemsep}{0pt}
\item $\mathbf{w}\in\operatorname{range}(T)$;
\item there exists $\psi\in\mathcal{W}'$ such that $\psi T=0$ and $\psi(\mathbf{w})\neq0$.
\end{enumerate}
\begin{proof}
\begin{enumerate}[label=(\roman*)]
\setlength{\itemsep}{0pt}
\item The alternatives are incompatible: if $\mathbf{w}=T\mathbf{v}$ and $\psi T=0$, then $\psi(\mathbf{w})=\psi(T\mathbf{v})=0$.
\item Suppose $\mathbf{w}\notin\operatorname{range}(T)$. Let $\mathcal{U}=\operatorname{range}(T)$, choose a basis $(\mathbf{u}_1,\ldots,\mathbf{u}_r)$ of $\mathcal{U}$, and extend $(\mathbf{u}_1,\ldots,\mathbf{u}_r,\mathbf{w})$ to a basis of $\mathcal{W}$. Define $\psi\in\mathcal{W}'$ by $\psi(\mathbf{u}_j)=0$, $\psi(\mathbf{w})=1$, and $\psi=0$ on the added basis vectors.\\
Then $\psi$ vanishes on $\operatorname{range}(T)$, so $\psi T=0$, while $\psi(\mathbf{w})=1$.
\end{enumerate}
\hlt{Structural Commentary}: A missing range vector is detected by a dual test that vanishes on the whole range.\\
\hlt{Key Insight}: Either a vector is reached by the map, or a functional separates it from the range.
\end{proof}
\end{theorem}

\begin{figure}[H]
  \centering
  \begin{tikzpicture}[
  >=Stealth,
  every node/.style={font=\small},
  ambient/.style={draw=gray!45, fill=gray!3, rounded corners=3pt, line width=0.55pt,
    minimum width=6.5cm, minimum height=3.35cm},
  subspace/.style={draw=teal!60!black, fill=teal!8, rounded corners=2pt,
    line width=0.55pt, minimum width=2.55cm, minimum height=0.72cm, align=center},
  point/.style={circle, fill=blue!70!black, inner sep=1.5pt},
  cert/.style={draw=red!60!black, dashed, line width=0.7pt},
  detect/.style={->, red!60!black, line width=0.7pt}
]
  \node[ambient] at (0,0) {};
  \node[anchor=north west] at (-3.02,1.52) {$\mathcal{W}$};

  \draw[cert] (-2.65,-0.40) -- (2.75,-0.40);
  \node[red!60!black, anchor=south west] at (-2.82,-0.34) {$\ker\psi$};

  \node[subspace] at (-0.55,-0.40) {$\operatorname{range}(T)$\\[-1pt]{\footnotesize $\psi=0$}};

  \node[point] (w) at (1.95,0.92) {};
  \node[anchor=west] at (2.08,0.92) {$\mathbf{w}$};
  \draw[detect] (w) -- (1.95,-0.32);
  \node[red!60!black, anchor=east] at (1.72,0.30) {$\psi(\mathbf{w})\neq0$};
\end{tikzpicture}

  \caption{A dual certificate detects a vector outside the range.}
  \label{fig:range-membership-dual-certificate}
\end{figure}

\begin{corollary} \label{cor:matrix-equation-left-nullspace-alternative} \hlt{Matrix Equation Alternative}\\
Let $A\in\operatorname{Mat}(m,n,\mathbb{F})$ and $\mathbf{b}\in\mathbb{F}^m$. Exactly one holds:
\begin{enumerate}[label=(\roman*)]
\setlength{\itemsep}{0pt}
\item $A\mathbf{x}=\mathbf{b}$ for some $\mathbf{x}\in\mathbb{F}^n$;
\item $\mathbf{y}^tA=\mathbf{0}^t$ and $\mathbf{y}^t\mathbf{b}\neq0$ for some $\mathbf{y}\in\mathbb{F}^m$.
\end{enumerate}
\begin{proof}
\begin{enumerate}[label=(\roman*)]
\setlength{\itemsep}{0pt}
\item By Theorem \ref{thm:range-membership-dual-alternative}, for $T_A:\mathbb{F}^n\to\mathbb{F}^m$ and $\mathbf{w}=\mathbf{b}$, $\mathbf{b}\in\operatorname{range}(T_A)$ is exactly $A\mathbf{x}=\mathbf{b}$ for some $\mathbf{x}\in\mathbb{F}^n$.
\item Under the standard dual identification, $\psi\in(\mathbb{F}^m)'$ has the form $\psi(\mathbf{z})=\mathbf{y}^t\mathbf{z}$.\\
Thus $\psi T_A=0$ iff $\mathbf{y}^tA=\mathbf{0}^t$, and $\psi(\mathbf{b})\neq0$ iff $\mathbf{y}^t\mathbf{b}\neq0$.
\end{enumerate}
\hlt{Structural Commentary}: Inconsistency is certified by a left-null vector.\\
\hlt{Key Insight}: Left null vectors are the coordinate form of dual separating functionals.
\end{proof}
\end{corollary}

\begin{corollary} \label{cor:surjective-iff-dual-injective} \hlt{Surjectivity and Dual Injectivity}\\
Let $\mathcal{V}$ and $\mathcal{W}$ be finite-dimensional, and let $T\in\LL(\mathcal{V},\mathcal{W})$. Then $T$ is surjective if and only if $T'$ is injective.
\begin{proof}
\begin{enumerate}[label=\roman*.]
\setlength{\itemsep}{0pt}
\item[$(\Rightarrow)$] Suppose $T$ is surjective. Then $\operatorname{range}(T)=\mathcal{W}$, so Theorem \ref{thm:dual-kernel-range-annihilator} gives $\ker(T')=\mathcal{W}^0=\{0\}$.
\item[$(\Leftarrow)$] Suppose $T'$ is injective. Then $\ker(T')=\{0\}$, so Theorem \ref{thm:dual-kernel-range-annihilator} gives $(\operatorname{range}(T))^0=\{0\}$. By Theorem \ref{thm:dimension-of-annihilator}, $\dim(\operatorname{range}(T))=\dim(\mathcal{W})$, hence $\operatorname{range}(T)=\mathcal{W}$.
\end{enumerate}
\hlt{Structural Commentary}: A map is onto exactly when no nonzero functional can vanish on its range. Dual injectivity detects this absence of invisible codomain directions.\\
\hlt{Key Insight}: Surjectivity of $T$ becomes injectivity of the pullback map $T'$.
\end{proof}
\end{corollary}

\begin{corollary} \label{cor:injective-iff-dual-surjective} \hlt{Injectivity and Dual Surjectivity}\\
Let $\mathcal{V}$ and $\mathcal{W}$ be finite-dimensional, and let $T\in\LL(\mathcal{V},\mathcal{W})$. Then $T$ is injective if and only if $T'$ is surjective.
\begin{proof}
\begin{enumerate}[label=\roman*.]
\setlength{\itemsep}{0pt}
\item[$(\Rightarrow)$] Suppose $T$ is injective. Then $\ker(T)=\{\mathbf{0}\}$, so Theorem \ref{thm:dual-kernel-range-annihilator} gives $\operatorname{range}(T')=(\ker(T))^0=\mathcal{V}'$.
\item[$(\Leftarrow)$] Suppose $T'$ is surjective. Then $\operatorname{range}(T')=\mathcal{V}'$, so Theorem \ref{thm:dual-kernel-range-annihilator} gives $(\ker(T))^0=\mathcal{V}'$. By Theorem \ref{thm:dimension-of-annihilator}, this forces $\dim(\ker(T))=0$, hence $\ker(T)=\{\mathbf{0}\}$.
\end{enumerate}
\hlt{Structural Commentary}: Injectivity says there are no nonzero domain directions killed by $T$. The dual statement says every test on the domain comes by pulling back a test from the codomain.\\
\hlt{Key Insight}: No kernel directions means no obstruction to extending functionals backward through $T$.
\end{proof}
\end{corollary}

\begin{theorem} \label{thm:matrix-of-dual-map-is-transpose} \hlt{Matrix of the Dual Map Is the Transpose}\\
Let $\mathcal{V}$ and $\mathcal{W}$ be finite-dimensional, let $T\in\LL(\mathcal{V},\mathcal{W})$, and let $\beta=(\mathbf{v}_1,\ldots,\mathbf{v}_n)$ and $\gamma=(\mathbf{w}_1,\ldots,\mathbf{w}_m)$ be bases of $\mathcal{V}$ and $\mathcal{W}$. If $\beta'$ and $\gamma'$ are the corresponding dual bases, then $[T']_{\beta'\leftarrow\gamma'}=\left([T]_{\gamma\leftarrow\beta}\right)^t$.
\begin{proof}
Let $A=[T]_{\gamma\leftarrow\beta}$ and $C=[T']_{\beta'\leftarrow\gamma'}$. Write $\gamma'=(\psi_1,\ldots,\psi_m)$ and $\beta'=(\varphi_1,\ldots,\varphi_n)$.\\
For $1\leq j\leq m$, $T'\psi_j=\sum\limits_{r=1}^n C_{r,j}\varphi_r$.
Evaluating at $\mathbf{v}_k$ gives $C_{k,j}=(T'\psi_j)(\mathbf{v}_k)=\psi_j(T\mathbf{v}_k)$.\\
Since $T\mathbf{v}_k=\sum\limits_{r=1}^m A_{r,k}\mathbf{w}_r$, we get $\psi_j(T\mathbf{v}_k)=A_{j,k}$. Hence $C_{k,j}=A_{j,k}$ for all $j,k$, so $C=A^t$.\\
\hlt{Structural Commentary}: Duality turns the matrix of a map into the matrix of its pullback on coordinate functionals. The row-column indexing is exactly transposition.\\
\hlt{Key Insight}: Coefficients of $T\mathbf{v}_k$ become values of the dual basis functionals on $T\mathbf{v}_k$.
\end{proof}
\end{theorem}

\begin{definition} \label{def:bilinear-form-and-matrix} \hlt{Bilinear Forms and Matrices}\\
A map $B:\mathcal{V}\times\mathcal{W}\to\mathbb{F}$ is \hlt{bilinear} if it is linear in each variable separately.\\
For bases $\beta=(\mathbf{v}_1,\ldots,\mathbf{v}_m)$ of $\mathcal{V}$ and $\gamma=(\mathbf{w}_1,\ldots,\mathbf{w}_n)$ of $\mathcal{W}$, the \hlt{matrix of $B$ in the bases $\beta,\gamma$} is $[B]_{\beta,\gamma}=[B(\mathbf{v}_i,\mathbf{w}_j)]$. If $\mathbf{x}=[\mathbf{v}]_{\beta}$ and $\mathbf{y}=[\mathbf{w}]_{\gamma}$, then $B(\mathbf{v},\mathbf{w})=\mathbf{x}^t [B]_{\beta,\gamma}\mathbf{y}$.
\end{definition}

\begin{proposition} \label{prop:bilinear-forms-linearize-through-duals} \hlt{Bilinear Forms Linearize through Duals}\\
Let $B:\mathcal{V}\times\mathcal{W}\to\mathbb{F}$ be bilinear, and define $\Phi_B:\mathcal{V}\to\mathcal{W}'$ where $(\Phi_B\mathbf{v})(\mathbf{w})=B(\mathbf{v},\mathbf{w})$,\\
and $\Psi_B:\mathcal{W}\to\mathcal{V}'$ where $(\Psi_B\mathbf{w})(\mathbf{v})=B(\mathbf{v},\mathbf{w})$. Then $\Phi_B,\Psi_B$ are linear.\\
For bases $\beta,\gamma$ and $M=[B]_{\beta,\gamma}$, $[\Phi_B]_{\gamma'\leftarrow\beta}=M^t$; and $[\Psi_B]_{\beta'\leftarrow\gamma}=M$.
\begin{proof}
For fixed $\mathbf{v}$ and $\mathbf{w}$, bilinearity gives $\Phi_B\mathbf{v}\in\mathcal{W}'$ and $\Psi_B\mathbf{w}\in\mathcal{V}'$. For $\mathbf{v}_1,\mathbf{v}_2\in\mathcal{V}$, $a\in\mathbb{F}$, and every $\mathbf{w}\in\mathcal{W}$, $(\Phi_B(a\mathbf{v}_1+\mathbf{v}_2))(\mathbf{w}) =B(a\mathbf{v}_1+\mathbf{v}_2,\mathbf{w}) =a(\Phi_B\mathbf{v}_1)(\mathbf{w})+(\Phi_B\mathbf{v}_2)(\mathbf{w})$, so $\Phi_B$ is linear.\\
For $\mathbf{w}_1,\mathbf{w}_2\in\mathcal{W}$, $a\in\mathbb{F}$, and every $\mathbf{v}\in\mathcal{V}$, $(\Psi_B(a\mathbf{w}_1+\mathbf{w}_2))(\mathbf{v}) =B(\mathbf{v},a\mathbf{w}_1+\mathbf{w}_2) =a(\Psi_B\mathbf{w}_1)(\mathbf{v})+(\Psi_B\mathbf{w}_2)(\mathbf{v})$,  so $\Psi_B$ is linear. Let $\beta=(\mathbf{v}_1,\ldots,\mathbf{v}_m)$ and $\gamma=(\mathbf{w}_1,\ldots,\mathbf{w}_n)$.\\
The $j$th $\gamma'$-coordinate of $\Phi_B(\mathbf{v}_i)$ is $(\Phi_B\mathbf{v}_i)(\mathbf{w}_j)=B(\mathbf{v}_i,\mathbf{w}_j)=M_{i,j}$, so column $i$ of $[\Phi_B]_{\gamma'\leftarrow\beta}$ is row $i$ of $M$ transposed, hence $[\Phi_B]_{\gamma'\leftarrow\beta}=M^t$. Likewise, the $i$th $\beta'$-coordinate of $\Psi_B(\mathbf{w}_j)$ is $M_{i,j}$, so $[\Psi_B]_{\beta'\leftarrow\gamma}=M$.\\
\hlt{Structural Commentary}: A bilinear form becomes linear once one input is frozen. Tracking both freezes gives two dual-valued linear maps whose matrices are transposes of one another.\\
\hlt{Key Insight}: Bilinearity is linearity into a dual space.
\end{proof}
\end{proposition}

\begin{perspective} \label{per:bilinear-forms-before-tensor-products} \hlt{Bilinear Forms before Tensor Products}\\
The maps $\Phi_B$ and $\Psi_B$ place bilinear forms inside dual-space and transpose machinery. The tensor model below is only a finite-coordinate tool; the full tensor product belongs later.
\end{perspective}

\begin{perspective} \label{per:subspace-sums-products-and-tensor-scaffold} \hlt{Sums, Products, and Tensor Scaffolds}\\
Direct sums and products separate coordinates; the tensor scaffold below separates paired basis directions as rank-one matrix slots.
\end{perspective}

\begin{figure}[H]
  \centering
  \begin{tikzpicture}[
  >=Stealth,
  every node/.style={font=\small},
  ambient/.style={draw=gray!45, fill=gray!3, rounded corners=3pt, line width=0.55pt,
    minimum width=4.55cm, minimum height=2.5cm},
  set/.style={draw=gray!55, fill=gray!6, line width=0.55pt},
  box/.style={draw=gray!45, fill=gray!3, rounded corners=3pt, line width=0.55pt,
    minimum width=2.45cm, minimum height=0.78cm, align=center},
  map/.style={->, gray!65!black, line width=0.75pt}
]
  \node[ambient] at (-2.15,0) {};
  \node[anchor=north west] at (-4.18,1.08) {$\mathcal{U}+\mathcal{W}$};

  \draw[set, fill=blue!7] (-2.75,0) ellipse (1.05cm and 0.65cm);
  \draw[set, fill=teal!8, fill opacity=0.75] (-1.55,0) ellipse (1.05cm and 0.65cm);
  \node at (-3.15,0.38) {$\mathcal{U}$};
  \node at (-1.14,0.38) {$\mathcal{W}$};
  \node at (-2.15,-0.02) {$\mathcal{U}\cap\mathcal{W}$};

  \node[box] (prod) at (2.25,0.55) {$\mathcal{U}\times\mathcal{W}$};
  \node[box] (sum) at (2.25,-0.75) {$\mathcal{U}+\mathcal{W}$};
  \draw[map] (prod) -- (sum);
  \node[gray!70!black, anchor=west] at (2.68,-0.10) {$\Gamma(\mathbf{u},\mathbf{w})=\mathbf{u}+\mathbf{w}$};
\end{tikzpicture}

  \caption{Intersection, sum, and product of subspaces.}
  \label{fig:subspace-sum-product-schematic}
\end{figure}

\begin{definition} \label{def:finite-coordinate-tensor-product-model} \hlt{Finite-Coordinate Tensor Product Model}\\
Let $\mathcal{V},\mathcal{W}$ have ordered bases $\beta=(\mathbf{v}_1,\ldots,\mathbf{v}_m)$ and $\gamma=(\mathbf{w}_1,\ldots,\mathbf{w}_n)$.\\
The \hlt{coordinate tensor model} for $(\beta,\gamma)$ is $\operatorname{Mat}(m,n,\mathbb{F})$, with $\mathbf{v}_i\otimes_{\beta,\gamma}\mathbf{w}_j$ represented by $E_{i,j}$.\\
For $\mathbf{v}=\sum_i x_i\mathbf{v}_i$ and $\mathbf{w}=\sum_j y_j\mathbf{w}_j$, set $\mathbf{v}\otimes_{\beta,\gamma}\mathbf{w} \defeq \sum\limits_{i=1}^m\sum\limits_{j=1}^n x_i y_j E_{i,j} \in\operatorname{Mat}(m,n,\mathbb{F})$.
\end{definition}

\begin{proposition} \label{prop:tensor-product-matrix-basis-model} \hlt{Tensor Model Basis and Bilinear Evaluation}\\
In the coordinate tensor model associated to bases $\beta=(\mathbf{v}_1,\ldots,\mathbf{v}_m)$ and $\gamma=(\mathbf{w}_1,\ldots,\mathbf{w}_n)$:
\begin{enumerate}[label=(\roman*)]
\setlength{\itemsep}{0pt}
\item the matrices $E_{i,j}=\mathbf{v}_i\otimes_{\beta,\gamma}\mathbf{w}_j$ form a basis;
\item the map $(\mathbf{v},\mathbf{w})\mapsto\mathbf{v}\otimes_{\beta,\gamma}\mathbf{w}$ is bilinear;
\item every bilinear form $B:\mathcal{V}\times\mathcal{W}\to\mathbb{F}$ determines a unique linear functional $\Theta_B$ on $\operatorname{Mat}(m,n,\mathbb{F})$ such that\\
$\Theta_B(\mathbf{v}\otimes_{\beta,\gamma}\mathbf{w})=B(\mathbf{v},\mathbf{w})$ for all $\mathbf{v}\in\mathcal{V}$ and $\mathbf{w}\in\mathcal{W}$.
\end{enumerate}
\begin{proof}
\begin{enumerate}[label=(\roman*)]
\setlength{\itemsep}{0pt}
\item This is the standard entry-matrix basis of $\operatorname{Mat}(m,n,\mathbb{F})$ from Proposition \ref{prop:matrix-spaces-are-vector-spaces}.
\item The coordinates of $a\mathbf{v}_1+\mathbf{v}_2$ are the same linear combination of the coordinates of $\mathbf{v}_1$ and $\mathbf{v}_2$, and the displayed formula is linear in those coordinates. The same argument applies in the second variable.
\item Define $\Theta_B(E_{i,j})=B(\mathbf{v}_i,\mathbf{w}_j)$ on the basis from (i), and extend linearly by Proposition \ref{prop:linear-map-determined-by-basis}. If $\mathbf{v}=\sum_i x_i\mathbf{v}_i$ and $\mathbf{w}=\sum_j y_j\mathbf{w}_j$, then $\Theta_B(\mathbf{v}\otimes_{\beta,\gamma}\mathbf{w}) =\sum\limits_{i=1}^m\sum\limits_{j=1}^n x_i y_j B(\mathbf{v}_i,\mathbf{w}_j) =B(\mathbf{v},\mathbf{w})$, by bilinearity of $B$.\\
Uniqueness follows since the values on basis matrices $E_{i,j}$ are forced by $\mathbf{v}_i,\mathbf{w}_j$.
\end{enumerate}
\hlt{Structural Commentary}: The coordinate tensor model stores each pair of basis directions in its own matrix slot, so a bilinear form can be evaluated by a single linear functional on those slots.\\
\hlt{Key Insight}: Tensor coordinates turn products of coordinates into linear coordinates.
\end{proof}
\end{proposition}

\begin{figure}[H]
  \centering
  \begin{tikzpicture}[
  >=Stealth,
  every node/.style={font=\small},
  map/.style={->, gray!65!black, line width=0.75pt},
  slot/.style={draw=gray!45, fill=gray!4, minimum width=0.68cm, minimum height=0.52cm},
  active/.style={slot, fill=teal!10, draw=teal!65!black, line width=0.65pt},
  basis/.style={draw=gray!45, fill=gray!3, rounded corners=3pt, line width=0.55pt,
    minimum width=2.45cm, minimum height=0.78cm, align=center}
]
  \node[basis] (pair) at (-2.45,0.52) {$\mathbf{v}_i\otimes_{\beta,\gamma}\mathbf{w}_j$};

  \draw[map] (pair.east) -- (-0.60,0.52);

  \foreach \x in {0,1,2,3} {
    \foreach \y in {0,1,2} {
      \node[slot] at (\x*0.70,\y*0.54) {};
    }
  }
  \node[active] at (1.40,0.54) {$E_{i,j}$};

  \draw[gray!50] (-0.02,0.54) -- (2.10,0.54);
  \draw[gray!50] (1.40,-0.26) -- (1.40,1.34);
\end{tikzpicture}

  \caption{A basis pair occupies the $(i,j)$ matrix slot.}
  \label{fig:tensor-product-matrix-basis}
\end{figure}

\begin{definition} \hlt{Double Dual Space}\\
The \hlt{double dual space} of $\mathcal{V}$ is $\mathcal{V}''\defeq(\mathcal{V}')'$. Define $\Lambda_{\mathcal{V}}:\mathcal{V}\to\mathcal{V}''$ by $(\Lambda_{\mathcal{V}}\mathbf{v})(\varphi)=\varphi(\mathbf{v}), \mathbf{v}\in\mathcal{V},\ \varphi\in\mathcal{V}'$.
\end{definition}

\begin{theorem} \label{thm:natural-map-to-double-dual} \hlt{Natural Map to the Double Dual}\\
The map $\Lambda_{\mathcal{V}}:\mathcal{V}\to\mathcal{V}''$ satisfies:
\begin{enumerate}[label=(\roman*)]
\setlength{\itemsep}{0pt}
\item $\Lambda_{\mathcal{V}}$ is linear;
\item if $T\in\LL(\mathcal{V},\mathcal{W})$ and $T''=(T')'$, then $T''\Lambda_{\mathcal{V}}=\Lambda_{\mathcal{W}}T$;
\item if $\mathcal{V}$ is finite-dimensional, then $\Lambda_{\mathcal{V}}$ is an isomorphism.
\end{enumerate}
\begin{proof}
\begin{enumerate}[label=(\roman*)]
\setlength{\itemsep}{0pt}
\item
For fixed $\mathbf{v}\in\mathcal{V}$, the map $\varphi\mapsto\varphi(\mathbf{v})$ is linear on $\mathcal{V}'$, so $\Lambda_{\mathcal{V}}\mathbf{v}\in\mathcal{V}''$.\\
For $\mathbf{u},\mathbf{v}\in\mathcal{V}$, $a\in\mathbb{F}$, and $\varphi\in\mathcal{V}'$, $(\Lambda_{\mathcal{V}}(a\mathbf{u}+\mathbf{v}))(\varphi) =\varphi(a\mathbf{u}+\mathbf{v}) =a(\Lambda_{\mathcal{V}}\mathbf{u})(\varphi)+(\Lambda_{\mathcal{V}}\mathbf{v})(\varphi)$, so $\Lambda_{\mathcal{V}}$ is linear.
\item If $\psi\in\mathcal{W}'$, then $\left(T''(\Lambda_{\mathcal{V}}\mathbf{v})\right)(\psi) =(\Lambda_{\mathcal{V}}\mathbf{v})(T'\psi) =(T'\psi)(\mathbf{v}) =\psi(T\mathbf{v}) =\left(\Lambda_{\mathcal{W}}(T\mathbf{v})\right)(\psi)$; $T''\Lambda_{\mathcal{V}}=\Lambda_{\mathcal{W}}T$.
\item
Suppose $\mathcal{V}$ is finite-dimensional. If $\mathbf{v}\neq\mathbf{0}$, extend $(\mathbf{v})$ to a basis of $\mathcal{V}$ and let $\varphi$ be the first functional of the corresponding dual basis. Then $\varphi(\mathbf{v})=1$, so $(\Lambda_{\mathcal{V}}\mathbf{v})(\varphi)=1$ and $\Lambda_{\mathcal{V}}\mathbf{v}\neq0$. Hence $\ker(\Lambda_{\mathcal{V}})=\{\mathbf{0}\}$ and $\Lambda_{\mathcal{V}}$ is injective. Since $\dim(\mathcal{V}'')=\dim(\mathcal{V}')=\dim(\mathcal{V})$, Corollary \ref{cor:equal-dimension-kernel-range} shows that $\Lambda_{\mathcal{V}}$ is an isomorphism.
\end{enumerate}
\hlt{Structural Commentary}: The double-dual map sends a vector to its evaluation action on functionals. It is natural because it commutes with every linear map.\\
\hlt{Key Insight}: Finite-dimensional vectors are completely recovered by how all linear functionals evaluate them.
\end{proof}
\end{theorem}

\begin{perspective} \hlt{Why the Double Dual Is Natural}\\
A finite-dimensional space $\mathcal{V}$ is isomorphic to $\mathcal{V}'$, but such an isomorphism usually requires choosing a basis. The map $\Lambda_{\mathcal{V}}:\mathcal{V}\to\mathcal{V}''$ requires no choice of basis.
\end{perspective}
